\section{Определения и обозначения}
\label{sec:Abbreviations} \index{Abbreviations}

\subsection{NUMA (Non-Uniform Memory Access)}

\textbf{NUMA} (Non-Uniform Memory Access)~--- архитектура многопроцессорных вычислительных систем с неоднородным доступом к оперативной памяти, при которой время доступа вычислительного узла к памяти зависит от её расположения относительно данного узла.

В \textbf{NUMA-системе} аппаратные ресурсы организованы в виде \textit{NUMA-нод}. Каждая нода включает один или несколько процессоров и привязанную к ним область оперативной памяти. Память NUMA-ноды называется \textit{локальной} для процессоров данной ноды и \textit{удалённой} для процессоров остальных NUMA-нод. Доступ к локальной памяти осуществляется быстрее, чем к удалённой.

Термин \textbf{NUMA-нода} используется для обозначения элементарной единицы топологии NUMA-системы, объединяющей вычислительные ресурсы и локальную память.

Под \textbf{NUMA-эффектом} понимается снижение производительности приложений, возникающее вследствие неоднородности доступа к памяти в NUMA-системах, в частности при обращении к удалённой памяти.

Под \textbf{cross-NUMA нагрузкой} понимается вычислительная нагрузка, выполнение которой охватывает несколько NUMA-нод и сопровождается активным межнодовым доступом к памяти. Такие нагрузки обычно не могут быть эффективно локализованы в пределах одной NUMA-ноды из-за объёма используемой памяти или характера взаимодействия потоков и данных.

\subsection{Таблицы трансляций}

Виртуальная память в современных системах реализуется на основе многоуровневой иерархии таблиц страниц. Основные уровни иерархии включают:
\begin{itemize}
    \item\textbf{PGD (Page Global Directory)}~--- верхний уровень иерархии.
    \item\textbf{PUD (Page Upper Directory)}~--- уровень, следующий за PGD.
    \item\textbf{PMD (Page Middle Directory)}~--- промежуточный уровень, содержащий ссылки на таблицы PTE или отображения \textit{больших страниц}.
    \item\textbf{PTE (Page Table Entry)}~--- конечный уровень, задающий отображения виртуальных страниц на физические.
\end{itemize}

\textbf{Записи} верхних уровней содержат ссылки на таблицы следующего уровня, тогда как \textit{записи} уровня PTE содержат физический адрес страницы и её атрибуты доступа.

Для ускорения трансляции адресов процессоры используют специальный аппаратный кэш --- \textbf{TLB} (\textit{Translation Lookaside Buffer}). В нём сохраняются результаты недавних преобразований виртуальных адресов в физические, что позволяет избежать повторного обхода иерархии таблиц трансляций при последующих обращениях к памяти. При изменении отображений в таблицах трансляций соответствующие записи TLB должны быть инвалидированы, чтобы процессор не использовал устаревшие результаты трансляции.

\subsection{Модели отображений и свойства памяти}

В контексте управления памятью в Linux используется несколько терминов, описывающих свойства памяти и её отображений.

\textbf{Private и shared отображения} определяются на уровне виртуальной памяти процесса (VMA). \textit{Private} отображение означает, что модификации страниц, выполненные одним процессом, не должны быть видны другим процессам, отображающим тот же участок памяти. Такие модификации реализуются посредством механизма \textit{copy-on-write}. \textit{Shared} отображение, напротив, предполагает, что изменения, выполненные одним процессом, становятся видны другим процессам.

\textbf{File-backed и анонимная память} описывают источник данных страниц памяти. \textit{File-backed} память отображает содержимое файлов или других разделяемых объектов, тогда как \textit{анонимная} память не связана с файловой системой и используется для динамически выделяемой памяти процесса.

\textbf{Read-only и writable память} определяют права доступа к страницам памяти. \textit{Read-only} страницы допускают только чтение (и, при необходимости, выполнение), тогда как \textit{writable} страницы допускают модификацию содержимого.

Отдельно следует отметить понятие \textbf{разделяемой памяти}. Несмотря на созвучность с термином \textit{shared}, данные понятия не являются эквивалентными. \textit{Shared} отображение определяется на уровне виртуального адресного пространства и отражает намерение процессов совместно использовать память. В то время как \textit{разделяемая} память~--- это состояние физической памяти, при котором одна и та же физическая страница используется в нескольких контекстах. В частности, физические страницы \textit{private} отображений также могут \textit{разделяться} между процессами до тех пор, пока один из процессов не попытается выполнить модификацию данных, что приведёт к срабатыванию механизма \textit{copy-on-write}.
